\documentclass[11pt]{article}

\usepackage[margin=1in]{geometry}
\usepackage{amsmath,amssymb,amsthm,mathtools,bm}
\usepackage{algorithm}
\usepackage{algorithmic}
\usepackage{booktabs,tabularx}
\usepackage{enumitem}
\usepackage{hyperref}
\usepackage[nameinlink]{cleveref}

\newtheorem{theorem}{Theorem}[section]
\newtheorem{lemma}[theorem]{Lemma}
\newtheorem{proposition}[theorem]{Proposition}
\newtheorem{corollary}[theorem]{Corollary}
\newtheorem{conjecture}[theorem]{Conjecture}
\theoremstyle{definition}
\newtheorem{definition}[theorem]{Definition}
\newtheorem{assumption}[theorem]{Assumption}

\newcommand{\norm}[1]{\left\lVert #1\right\rVert}
\newcommand{\vol}{\operatorname{vol}}
\newcommand{\supp}{\operatorname{supp}}
\newcommand{\eps}{\epsilon}
\newcommand{\softO}{\widetilde{\mathcal O}}
\newcommand{\AESP}{\textsc{AESP}}
\newcommand{\LOCSOR}{\textsc{LocSOR}}
\newcommand{\APPR}{\textsc{APPR}}

\title{Hybrid Local Solvers for Personalized PageRank\\
Independent Research Note}
\author{Baojian Zhou}
\date{August 2026}

\begin{document}
\maketitle

\begin{abstract}
This note organizes the hybrid local solver research program into an
independent mathematical document. The central idea is to combine accelerated
outer convergence from \AESP{} with a monotone local relaxation tail. The note
separates established results, conditional complexity bounds, and unresolved
locality questions.
\end{abstract}

\tableofcontents

\section{Scope and research status}
This note is a rigorous synthesis of the hybrid local solver research discussion.

The goal is to develop local PageRank solvers with accelerated dependence on
$\alpha$. Classical local push methods achieve worst-case complexity
$\mathcal O(1/(\alpha\eps))$, while accelerated methods suggest a target scale
$\softO(1/(\sqrt\alpha\eps))$.

The note follows a strict claim status convention:

\begin{itemize}
\item source results: directly inherited from published algorithms;
\item proved results: derived in this note;
\item conditional results: requiring explicit locality assumptions;
\item open problems: not currently established.
\end{itemize}

The main unresolved issue is not convergence of the hybrid method, but proving
that accelerated burn-in preserves sufficiently small local work.


\section{Problem formulation}
% Canonical source-aligned PageRank/RPPR reference formulation.
%
% This fragment is intentionally heading-free at the section level so that the
% active paper and every standalone note can input it. It is a manuscript
% reference convention, not an implementation-wide residual decision.
% Primary source: Fountoulakis and Martinez-Rubio, arXiv:2602.21138v2,
% PDF p. 3, Section 3 and equation (RPPR).

\subsection{Common source-aligned PageRank and RPPR model}
\label{subsec:shared-source-aligned-problem}

All documents in this manuscript workspace use the following reference
language. It follows the plain italic notation of the comparison paper:
\(x,s,A,D,Q,I\) denote column vectors or matrices as determined by their
displayed domains. This is the scoped exception to the project's usual
bold-vector and bold-matrix typography. A note may impose a stronger
assumption after this block, but it must not redefine these objects.

Let \(G=(V,E)\) be a finite, simple, undirected, unweighted graph without
isolated vertices, let \(V=[n]:=\{1,\ldots,n\}\), and let
\(A\in\{0,1\}^{n\times n}\) be its symmetric adjacency matrix. For
\(i\in V\), write
\[
    \mathcal N(i):=\{j\in V:\{i,j\}\in E\},
    \qquad
    d_i:=|\mathcal N(i)|,
    \qquad
    D:=\diag(d_1,\ldots,d_n).
\]
For \(S\subseteq V\), define
\[
    \mathcal N(S):=\bigcup_{i\in S}\mathcal N(i),
    \qquad
    \partial S:=\mathcal N(S)\setminus S,
    \qquad
    \operatorname{Ext}(S):=V\setminus(S\cup\partial S),
\]
\begin{equation}
    \vol(S):=\sum_{i\in S}d_i.
    \label{eq:shared-volume}
\end{equation}
The support of a vector is
\(\supp(x):=\{i\in[n]:x_i\neq0\}\). Matrix inequalities use the Loewner
order, and all unqualified vector inequalities are coordinatewise.

The seed is a column distribution
\begin{equation}
    s\in\R_+^n,
    \qquad
    \inner{\one}{s}=1.
    \label{eq:shared-seed}
\end{equation}
The single-seed specialization is always written \(s=e_v\), where \(v\in V\)
is the seed vertex. Thus \(s\) is never used as a vertex label. For
\(\alpha\in(0,1]\), define
\begin{equation}
\begin{aligned}
    \mathcal L&:=I-D^{-1/2}AD^{-1/2},\\
    Q&:=\alpha I+\frac{1-\alpha}{2}\mathcal L
      =\frac{1+\alpha}{2}I
       -\frac{1-\alpha}{2}D^{-1/2}AD^{-1/2},\\
    b&:=\alpha D^{-1/2}s.
\end{aligned}
    \label{eq:shared-pagerank-matrices}
\end{equation}
Since \(0\preceq\mathcal L\preceq2I\),
\(\alpha I\preceq Q\preceq I\). The shared smooth PageRank quadratic is
\begin{equation}
    f(x):=\frac12\inner{x}{Qx}-\inner{b}{x},
    \qquad
    \nabla f(x)=Qx-b.
    \label{eq:shared-pagerank-objective}
\end{equation}
Its unique unregularized minimizer and associated lazy personalized PageRank
vector are
\begin{equation}
    x^0:=Q^{-1}b,
    \qquad
    \pi:=D^{1/2}x^0.
    \label{eq:shared-ppr-solution}
\end{equation}

For a regularization parameter \(\rho>0\), reserve \(g_\rho\) for the
nonsmooth weighted \(\ell_1\) term and define
\begin{equation}
    g_\rho(x):=\alpha\rho\norm{D^{1/2}x}_1,
    \qquad
    F_\rho(x):=f(x)+g_\rho(x),
    \qquad
    x^\star(\rho):=\argmin_{x\in\R^n}F_\rho(x).
    \label{eq:shared-rppr-objective}
\end{equation}
When \(\rho\) is fixed, \(x^\star\) abbreviates \(x^\star(\rho)\), never the
unregularized point \(x^0\). Write
\[
    S^\star(\rho):=\supp(x^\star(\rho)),
    \qquad
    I^\star(\rho):=[n]\setminus S^\star(\rho).
\]
The source records \(x^\star(\rho)\geq0\),
\(\vol(S^\star(\rho))\leq1/\rho\), and the coordinatewise conditions
\begin{equation}
    \nabla_i f(x^\star(\rho))
    \in
    \begin{cases}
        \{-\alpha\rho\sqrt{d_i}\}, & x_i^\star(\rho)>0,\\
        [-\alpha\rho\sqrt{d_i},0], & x_i^\star(\rho)=0.
    \end{cases}
    \label{eq:shared-rppr-kkt}
\end{equation}

The common computational unit is one repeated adjacency-list scan: scanning
coordinate \(i\) costs \(d_i\), and scanning a set \(S\) costs \(\vol(S)\).
Iteration count and wall-clock time do not replace this degree-weighted work
measure.

Finally, accuracy symbols remain distinct. The APPR threshold is
\(\epsappr\); \(\epsobj\) denotes an objective-gap target;
\(\epspg\) denotes the source experiment's proximal fixed-point tolerance; and
\(\epsppr\) may denote a document-scoped degree-normalized PPR error target.
No equality or conversion among these quantities is assumed. In particular,
this shared model does not resolve the repository-wide residual decision.


\paragraph{Note specialization.}
This note assumes \(0<\alpha<1/2\). Define the proof-scoped degree vector and
weighted mass functional
\[
  w:=D^{1/2}\one,
  \qquad
  p(x):=w^Tx.
\]
The identities
\begin{equation}
  Qw=\alpha w,
  \qquad
  w^TQ=\alpha w^T,
  \qquad
  w^Tb=\alpha
  \label{eq:weighted-identities}
\end{equation}
are used throughout.  The matrix $Q$ is a symmetric positive-definite
$M$-matrix and satisfies $\alpha I\preceq Q\preceq I$.

\subsection{Orthant form and note-scoped KKT residual}

Write \(J(x):=\norm{D^{1/2}x}_1=\sum_i w_i|x_i|\). Since \(b\ge0\)
and \(Q\) is an \(M\)-matrix, the shared minimizer \(x^\star(\rho)\) is
nonnegative. In this note it is convenient to write it as \(x_\rho^\star\)
and use the equivalent orthant formulation
\begin{equation}
 F_\rho^+(x)
 :=\frac12x^TQx-b^Tx+\alpha\rho w^Tx+I_{\R_+^V}(x).
 \label{eq:orthant-rppr}
\end{equation}

On the orthant, introduce the note-scoped signed KKT residual
\begin{equation}
  \kappa_\rho(x):=b-Qx-\alpha\rho w.
  \label{eq:kkt-residual}
\end{equation}
The minimizer obeys
\begin{equation}
  \kappa_\rho(x_\rho^\star)\le0,
  \qquad
  x_\rho^\star\ge0,
  \qquad
  x_{\rho,i}^\star \kappa_{\rho,i}(x_\rho^\star)=0.
  \label{eq:kkt}
\end{equation}
Our sign convention makes positive residual a legal direction in which to
increase a coordinate.  This convention is local to this note.

Multiplying the active KKT equations by $w$ gives the standard support budget
\begin{equation}
  p(x_\rho^\star)+\rho\vol(\supp x_\rho^\star)\le1.
  \label{eq:rppr-support-budget}
\end{equation}
Thus $\rho^{-1}$ is simultaneously a regularization scale and a worst-case
support-volume scale. It is not identified here with an unresolved global
accuracy symbol.

\subsection{Local work}

Scanning the neighborhood of vertex $i$ costs $d_i$. By the shared work
definition, a batched active set $S$ therefore costs $\vol(S)$.
For an active-set trajectory $(S_t)_{t<T}$ we charge
\begin{equation}
  \Work_T:=\sum_{t=0}^{T-1}\vol(S_t).
  \label{eq:work-model}
\end{equation}
This is the same edge-operation unit used in the locally evolving-set
framework of \citet{zhou2024iterative}, although different algorithms may
choose different active sets.

The target~\eqref{eq:target-work} is stronger than multiplying an
$O(1/\sqrt\alpha)$ iteration count by the trivial $O(1/\rho)$ support bound
only if each accelerated subproblem can actually be solved with local work and
the momentum state itself does not create expensive transient support.


\section{Source frameworks}
\subsection{Three relevant baselines}

\paragraph{Monotone local push.}
The approximate PageRank push method of \citet{andersen2007using} is local
because nonnegative residual mass is spent monotonically.  In its native
degree-normalized tolerance, its classical worst-case work is
$O(1/(\alpha\epsappr))$.  The repository contains a separate
matching star lower bound; this note uses the result only as motivation and
does not assert a conversion between $\epsappr$ and $\rho$.

\paragraph{Locally evolving sets.}
\citet{zhou2024iterative} define a local process with
$S_{t+1}\subseteq S_t\cup N(S_t)$ and charge~\eqref{eq:work-model}.  Their
analysis emphasizes that nonmonotone active sets can still be local when
active volume is paired with a residual-progress factor.  This is the natural
work model for the present discussion: convergence in a global norm alone is
not enough.

\paragraph{Accelerated outer processes.}
Catalyst approximately solves a sequence of quadratically regularized
subproblems and extrapolates their centers
\citep{lin2018catalyst}.  AESP localizes this idea for PageRank and uses only
$\softO(1/\sqrt\alpha)$ outer subproblems
\citep{huang2025accelerated}.  AESP's published work contains a run-dependent
locality quantity and its experiments show especially rapid progress in the
early stages.  Classical FISTA, however, can activate a high-degree star
center that ISTA never touches \citep{fountoulakis2026complexity}.  Acceleration
and locality must therefore be analyzed jointly.

\subsection{The bell-volume observation}

The motivating empirical trajectory has three phases:
\begin{enumerate}[label=(\roman*),leftmargin=*]
  \item a small active set and rapid accelerated error reduction;
  \item a middle phase in which active volume expands and may form a large
  bell-shaped peak; and
  \item a late phase in which a monotone local relaxation can finish cheaply
  once the remaining positive residual mass is small.
\end{enumerate}
This suggested two related questions.
\begin{itemize}[leftmargin=*]
  \item Can one stop acceleration before the volume peak and hand the state to
  a monotone local solver?
  \item Can the bell be flattened by allowing
  $\softO((1/\sqrt\alpha)\log(1/\rho))$ outer stages while enforcing
  $O(1/\rho)$ active volume at every stage?
\end{itemize}
The second question is not answered merely by the optimal support
budget~\eqref{eq:rppr-support-budget}: extrapolated intermediate points need
not share the optimum's support.

\subsection{Why the known sparse accelerated algorithm does not settle the target}

The expanding-subspace method of \citet{martinezrubio2023accelerated} gives a
rigorous accelerated sparse algorithm for positive-definite $M$-matrix
quadratics.  Its dependence on the optimal support differs from the desired
linear $1/\rho$ locality scale.  It proves that sparse acceleration is
possible, but it does not imply~\eqref{eq:target-work} for the classical
one-gradient or one-proximal-center trajectory studied here.


\section{Hybrid algorithm}
\subsection{Over-regularized warmup}

Set
\begin{equation}
  \bar\rho:=(1+\sqrt\alpha)\rho.
  \label{eq:overrho}
\end{equation}
The proof-oriented hybrid template is:
\begin{enumerate}[label=\arabic*.,leftmargin=*]
  \item run an accelerated local method on $F_{\bar\rho}$ until a certified
  coarse objective gap is reached;
  \item retract the accelerated point to a lower certificate
  $\ell\le x_{\bar\rho}^\star$ and discard the momentum state; and
  \item run a one-sided $\omega=1$ local Gauss--Seidel/proximal coordinate
  tail on $F_\rho$.
\end{enumerate}
The separation between the accelerated point and the retracted checkpoint is
essential.  Momentum is not itself a valid warm-start invariant for a
monotone residual process.

The choice~\eqref{eq:overrho} creates KKT slack of scale
$\alpha(\bar\rho-\rho)=\alpha^{3/2}\rho$.  This is large enough to make the
tail cheap, but changes the support budget only by a $1+O(\sqrt\alpha)$
factor.

\subsection{Why the tail uses \texorpdfstring{$\omega=1$}{omega=1}}

For a quadratic coordinate update, the exact objective change is negative for
every fixed relaxation $0<\omega<2$.  Objective descent alone is not the
locality invariant needed here: when $\omega>1$, signed residuals may be
created and must be charged separately.  The signed warm-start SOR lemma used
in the broader project allows $\omega\in[1,2]$ under its own hypotheses, but
the unconditional one-sided mass argument below uses the safe endpoint
$\omega=1$.  The note therefore makes no claim that the optimized SOR
parameter preserves the same tail bound.

\subsection{The desired division of work}

Let $T=\softO(1/\sqrt\alpha)$ denote the number of accelerated proximal
centers needed for the warmup gap.  The intended theorem decomposes as
\[
  W_{\mathrm{hybrid}}
  =W_{\mathrm{warm}}+W_{\mathrm{tail}}.
\]
The tail term is already controlled at the target scale.  Every unsuccessful
proof in this note fails at the same place: bounding $W_{\mathrm{warm}}$
without silently charging the whole graph or assuming monotone accelerated
support.


\section{Convergence and complexity}
For the local SOR tail, a coordinate update satisfies an exact quadratic descent
identity. Therefore any fixed relaxation parameter $0<\omega<2$ gives objective
descent.

With the safe choice $\omega=1$, an objective-gap handoff
\[
f(x_J)-f^\star\leq\frac{\alpha^{3/2}\eps}{1+\alpha}
\]
gives tail work
\[
\mathcal O\left(\frac1{\sqrt\alpha\eps}\right).
\]

The full hybrid complexity remains conditional on controlling AESP burn-in
locality, for example through
\[
\Lambda_J=\max_t\frac{\overline{\vol}(S_t)}{\gamma_t}=O(1/\eps).
\]


\section{Open problems}
\subsection{What is settled}

The discussion yields the following durable conclusions.
\begin{enumerate}[label=\arabic*.,leftmargin=*]
  \item The over-regularized handoff and the one-sided $\omega=1$ local tail
  are complete in the note-scoped KKT convention; the tail costs
  $O(1/(\rho\sqrt\alpha))$.
  \item A matched activation-tax proximal subproblem is individually local and
  can be solved locally at a constant inner contraction rate.
  \item The same linear activation tax prevents uniform exponent-two Bregman
  acceleration across changing support.  The corresponding dual envelope is
  not globally convex.
  \item A Euclidean proximal subproblem centered at a lower certificate is
  sandwiched below the RPPR optimum and has support volume at most $1/\rho$.
  \item For classical inertial Euclidean centers, cumulative active volume is
  controlled exactly by~\eqref{eq:volume-via-flux}.  Support growth itself is
  not the obstruction; repeated mass deactivation is.
  \item Branching trees support outward activation/deactivation waves.  The
  tree mechanism rules out using an unguarded $O(\alpha T)$ flux assumption in
  a graph-uniform proof.
  \item The product lower bound
  $\Omega(1/(\rho\sqrt\alpha))$ is rigorous for persistent-support one-hop
  methods, but not for arbitrary moving-frontier local algorithms.
\end{enumerate}

Accordingly, the unconditional graph-uniform work bound
\begin{equation}
  \Work_{\mathrm{hybrid}}
  =\softO\!\left(\frac1{\rho\sqrt\alpha}\right)
  \label{eq:still-open-target}
\end{equation}
must remain open for the classical extrapolated warmup.  It should not be
promoted into the active manuscript.

\subsection{The most credible safeguarded algorithm}

The tree calculation suggests a concrete redesign rather than another global
potential.  Maintain two states:
\begin{itemize}[leftmargin=*]
  \item an accelerated trial state, allowed to extrapolate; and
  \item a monotone lower shadow $\ell_t$, updated only by certified positive
  residual coordinate steps.
\end{itemize}
During an acceleration epoch, monitor a local wave certificate such as
\begin{equation}
 \mathfrak D_{a:b}
 :=\pos{\beta_{\mathrm A} p_{b-1}-p_b}
   +\beta_{\mathrm A}\sum_{t=a}^b
     p_{A_{t-2}\setminus A_t}(x_{t-2}),
 \label{eq:epoch-flux}
\end{equation}
or the weighted negative part suggested by~\eqref{eq:transformed-wave}.
When the certificate exceeds a fixed fraction of the source mass injected in
the epoch, discard momentum, retract to the lower shadow, enlarge the certified
face if necessary, and restart.

The observable theorem to seek is not that restarts never occur.  It is an
amortized statement of the form
\begin{equation}
  \Work_{\mathrm{total}}
  =\softO\!\left(
      \frac1{\rho\sqrt\alpha}
      +\frac{N_{\mathrm{rst}}}{\rho}
    \right),
  \label{eq:restart-target}
\end{equation}
where $N_{\mathrm{rst}}$ is the number of certified wave restarts, followed by
either a structural bound on $N_{\mathrm{rst}}$ or an instance-adaptive
guarantee that exposes it in the theorem. This formulation is compatible with
\cref{thm:bounded-core,thm:controlled-faces} and with classical adaptive
restart ideas \citep{odonoghue2015adaptive}, while using a graph-local signal
instead of a global objective test.

\subsection{Central missing lemmas}

\begin{description}[style=nextline,leftmargin=3.4cm]
  \item[Finite-tree excitation.] Complete the rigorous transfer from the
  outward period-11 wave to the zero-start root-source response on a sequence
  of finite simple trees.  This would convert the current computational and
  limiting evidence into a formal asymptotic counterexample to
  \cref{conj:flux}.

  \item[Restart amortization.] Prove that a wave-triggered restart pays for
  itself through objective decrease, newly certified KKT slack, or expansion
  of a bounded core.  Counting restarts without such a payment simply
  reintroduces the unknown $K$ in~\eqref{eq:face-phase-work}.

  \item[Lower-shadow accuracy.] Bound the gap between the accelerated trial
  point and the monotone lower shadow tightly enough that restarting does not
  lose the $1/\sqrt\alpha$ progress already obtained.

  \item[Flattened bell.] Determine whether
  $\softO((1/\sqrt\alpha)\log(1/\rho))$ guarded epochs can enforce
  $O(1/\rho)$ active volume per epoch.  The tree wave shows that a simple
  fixed momentum schedule cannot be the entire proof.

  \item[Broader oracle lower bound.] Either construct a hard family that
  forces simultaneous volume on a moving-frontier evolving-set algorithm, or
  prove that such an algorithm can evade the product lower bound.  Output size
  and iteration complexity alone do not settle this question.

  \item[Residual reconciliation.] Translate the note-scoped KKT threshold,
  the AESP degree-normalized infinity error, the APPR threshold, and the
  proximal fixed-point residual into one accepted repository convention before
  promoting any theorem.
\end{description}

\subsection{Final research position}

The hybrid idea remains viable, but the proof target has changed.  The
correct algorithm cannot be ``classical acceleration plus locality for free.''
It must make locality an enforced invariant or an observable restart
condition.  The strongest current theorem is the flux-dependent
bound~\eqref{eq:hybrid-flux-bound}, together with the bounded-core and
controlled-face corollaries.  The strongest obstruction is the combination
of the exact face-change defect~\eqref{eq:face-change-defect} and the rooted
tree wave~\eqref{eq:radial-obstacle}.  Any next proof should state explicitly
which of these two mechanisms it suppresses.


\end{document}
